\chapter{Integrable RG Flows and Perturbed Two-Dimensional CFT}
\label{ch:integrable-rg-flows-perturbed-cft}

\ConventionsMixedCFT

This chapter studies massive and massless two-dimensional flows that are
specified in the ultraviolet by a conformal field theory and a relevant local
operator, and whose infrared dynamics is constrained by factorized scattering.
The comparison data along the flow are finite-size ground-state energies,
local correlators, and scattering data, all expressed in declared
normalization conventions.

\section{Relevant Perturbations}

Let \(\mathcal C\) be a unitary Euclidean two-dimensional CFT on the plane
with local scalar primary operator \(\mathcal O(z,\bar z)\) of scaling
dimension
\[
  \Delta_{\mathcal O}=h+\bar h<2 .
\]
A regulated perturbation of \(\mathcal C\) is a family of correlation
functionals
\[
  \left\langle X\right\rangle_{g,\Lambda}
  =
  \frac{
  \left\langle
  X\exp\!\left[-g_\Lambda\int d^2x\,\mathcal O_\Lambda(x)\right]
  \right\rangle_{\mathcal C,\Lambda}}
  {\left\langle
  \exp\!\left[-g_\Lambda\int d^2x\,\mathcal O_\Lambda(x)\right]
  \right\rangle_{\mathcal C,\Lambda}},
\]
where \(\Lambda\) is a short-distance regulator, \(g_\Lambda\) is a
cutoff-dependent coordinate, \(\mathcal O_\Lambda\) is the corresponding
regularized insertion, and \(X\) is a product of separated local insertions.
The renormalized coupling has engineering dimension \(2-\Delta_{\mathcal O}\).
When the perturbation produces a massive theory with mass scale \(M\), the
dimensionless ratio
\[
  r=MR
\]
is the natural variable for placing the theory on a spatial circle of
circumference \(R\).

\section{Finite-Size Scaling Function}

Let \(E_0(R)\) be the ground-state energy on a circle.  Define the effective
central-charge function
\[
  c_{\mathrm{eff}}(R)
  =
  -\frac{6R}{\pi}E_0(R).
\]
For a CFT with periodic boundary conditions,
\[
  E_0(R)=-\frac{\pi c}{6R},
\]
so \(c_{\mathrm{eff}}(0)\) equals the ultraviolet central charge when the
scaling limit exists and no zero-mode obstruction changes the cylinder
vacuum.  For a massive infrared theory with a unique vacuum, the large-\(R\)
energy has the form
\[
  E_0(R)=\mathcal E_{\mathrm{bulk}}R+o(R^0),
\]
after the universal Casimir term disappears.  The bulk energy density
\(\mathcal E_{\mathrm{bulk}}\) is a local counterterm coordinate; comparisons
of \(c_{\mathrm{eff}}\) between ultraviolet and infrared theories must specify
whether this extensive term has been subtracted.

\section{TBA Scaling Function}

For a diagonal factorized massive theory, Chapter~\ref{ch:integrable-thermodynamic-bethe-ansatz}
gives pseudoenergies \(\epsilon_a(\theta;r)\) satisfying
\[
  \epsilon_a(\theta;r)
  =
  r\,\mu_a\cosh\theta
  -\sum_b\int_{\mathbb R}d\theta'\,
  \varphi_{ab}(\theta-\theta')
  \log\!\left(1+e^{-\epsilon_b(\theta';r)}\right),
\]
where \(m_a=\mu_a M\).  The finite-size ground-state energy in the
mirror channel is
\[
  E_0(R)
  =
  -\frac{M}{2\pi}
  \sum_a\mu_a\int_{\mathbb R}d\theta\,
  \cosh\theta\,
  \log\!\left(1+e^{-\epsilon_a(\theta;r)}\right)
  +\mathcal E_{\mathrm{bulk}}R .
\]
Consequently the subtracted scaling function is
\[
  c_{\mathrm{TBA}}(r)
  =
  \frac{3r}{\pi^2}
  \sum_a\mu_a\int_{\mathbb R}d\theta\,
  \cosh\theta\,
  \log\!\left(1+e^{-\epsilon_a(\theta;r)}\right).
\]
This formula is a construction from scattering data and an assumption that the
mirror-channel TBA computes the finite-size vacuum energy.

\begin{proposition}[Dilogarithm form of the constant-kernel ultraviolet limit]
\label{prop:tba-dilogarithm-uv}
Assume that as \(r\to0\) the functions \(\epsilon_a(\theta;r)\) have a plateau
region in which the kernels may be replaced by their integrals
\[
  N_{ab}=\int_{\mathbb R}d\theta\,\varphi_{ab}(\theta),
\]
and that the plateau values \(Y_a=e^{\epsilon_a}\) solve
\[
  \log Y_a=-\sum_b N_{ab}\log(1+Y_b^{-1}) .
\]
Then the plateau contribution to the ultraviolet central charge is
\[
  c_{\mathrm{UV}}
  =
  \frac{6}{\pi^2}\sum_a
  L\!\left(\frac{1}{1+Y_a}\right),
\]
where \(L(x)\) is the Rogers dilogarithm.
\end{proposition}

\begin{proof}
In the plateau approximation, the derivative of the TBA equation with respect
to \(\theta\) relates \(d\epsilon_a\) to the derivative of the driving term.
After integrating by parts in the scaling function, the contribution from
each species takes the form
\[
  \frac{6}{\pi^2}
  \int_{\epsilon_a(-\infty)}^{\epsilon_a(+\infty)}
  d\epsilon\,
  \log(1+e^{-\epsilon})
  \frac{d}{d\epsilon}\log(1+e^{\epsilon})^{-1}.
\]
With \(x=(1+e^\epsilon)^{-1}\), this integral is the Rogers dilogarithm
\[
  L(x)
  =
  -\frac12\int_0^x dt\,
  \left(
  \frac{\log(1-t)}{t}
  +\frac{\log t}{1-t}
  \right),
\]
evaluated at the plateau value.  The assumed algebraic equations determine
those endpoint values.  The stated formula follows by summing over species.
\end{proof}

\section{Massless Flows}

For an integrable flow between two CFTs, the infrared may contain massless
right- and left-moving particles.  A right mover has
\[
  E=p=\frac{M}{2}e^\theta ,
\]
and a left mover has
\[
  E=-p=\frac{M}{2}e^{-\theta}.
\]
The TBA then uses pseudoenergies
\(\epsilon_R(\theta)\) and \(\epsilon_L(\theta)\), together with
right-right, left-left, and right-left kernels.  The mixed kernel records the
interaction between the two chiral sectors along the flow.  The ultraviolet
and infrared central charges are extracted from the small-\(r\) and large-\(r\)
limits of the same finite-size functional after the bulk term is fixed.

\section{Operator Data Along the Flow}

Scattering data and TBA do not determine all local correlators by themselves.
To compare the perturbing CFT with the massive theory, one must also specify
form factors of local operators and their ultraviolet normalization.  For a
local operator \(\widehat{\mathcal O}\), the short-distance scaling dimension
is tested by the two-point function
\[
  \left\langle \widehat{\mathcal O}(x)\widehat{\mathcal O}(0)\right\rangle
  \sim
  \frac{C_{\mathcal O}}{|x|^{2\Delta_{\mathcal O}}}
  \quad (M|x|\to0),
\]
with \(C_{\mathcal O}\) fixed by a source convention.  A factorized model
therefore represents an integrable RG flow only after the scattering theory,
TBA scaling function, and local form-factor normalization agree with a common
ultraviolet operator chart.

\begin{openproblem}[Rigorous integrable-flow reconstruction]
\label{op:integrable-flow-reconstruction}
Construct a theorem that starts from factorized scattering data, form-factor
solutions satisfying locality and convergence estimates, and TBA finite-size
data, and reconstructs a local two-dimensional QFT whose ultraviolet limit is
a specified CFT perturbation.
\end{openproblem}
